\chapter{Prix juste}\label{ch:hf:fair-value}

Une même action est cotée sur une douzaine de plateformes, son contrat à terme sur indice se négocie à Chicago et le fonds indiciel coté (ETF) qui la détient se négocie dans le New Jersey. Le
prix du teneur de marché n'est aucune de ces cotations mais sa propre estimation, révisée à chaque message en provenance de l'une quelconque d'entre elles. Dans le marché simulé de ce chapitre,
où un instrument suiveur se négocie sur deux plateformes et un meneur bouge en premier, la meilleure estimation construite à partir du carnet propre à l'instrument suiveur prévoit
son mid cinq secondes à l'avance 2.4~\% mieux que le mid ne le fait. Ajouter le prix du meneur améliore la prévision de 3.5~\% lorsque le meneur bouge une demi-
seconde avant, et de 12.8~\% lorsqu'il bouge deux secondes avant. De faibles pourcentages, appliqués à chaque cotation de chaque journée, sont l'objet de ce chapitre.

\section{Mid, mid pondéré et microprix}

\begin{definition}[Prix juste]\label{def:hf:fair-value:fair}
Le \emph{prix juste}\index{prix juste} (fair price) $\hat P_t$ d'un teneur de marché est son estimation, à l'instant $t$, du prix efficient $P^\ast_t$ de l'instrument
à partir de l'information dont il dispose à $t$~: les carnets des plateformes où l'instrument se négocie, les prix des instruments liés et ses exécutions récentes.
Les cotations se placent autour de lui~; l'écart entre le prix d'une exécution et le prix juste au moment de celle-ci est le profit espéré de cette exécution.
\end{definition}

Le nom est choisi par rapport à deux notions voisines. La \emph{fair value} d'un contrat à terme (livre~1, chapitre~21) est un prix sans arbitrage à partir du spot,
du financement et des dividendes~; la \emph{valeur théorique} d'une option (livre~5, chapitre~26) provient d'une surface de volatilité. Un \omterm{def:hf:fair-value:fair}{prix juste} peut
utiliser l'un ou l'autre comme donnée d'entrée, mais c'est une estimation statistique, révisée à chaque message, et sa qualité se mesure à la précision de ses prévisions.

Trois estimateurs n'exploitent que le meilleur niveau du carnet d'une plateforme. Le mid $m_t=\tfrac12(b_t+a_t)$ ignore les tailles des files $Q^b_t$ et $Q^a_t$. Le mid pondéré
\[
w_t=\frac{a_t\,Q^b_t+b_t\,Q^a_t}{Q^b_t+Q^a_t}
\]
se rapproche du cours vendeur lorsque la file acheteuse est la plus grande, car une grande file acheteuse a moins de chances d'être épuisée en premier. Le microprix (livre~7,
chapitre~8) remplace la formule fixe par une table apprise sur les données~: le mid plus la variation moyenne du mid sur un horizon, par
tranche de déséquilibre des files $I_t=Q^b_t/(Q^b_t+Q^a_t)$. Stoikov a montré que le microprix ainsi estimé prévoit mieux les prix à court terme que le
mid ou le mid pondéré.

\begin{example}[Trois estimations d'un même carnet]\label{ex:hf:fair-value:three}
Un carnet affiche 300 actions au cours acheteur de 99 ticks et 100 au cours vendeur de 100. Le mid vaut 99.5~; le déséquilibre vaut 0.75 et le mid pondéré 99.75. La
table de microprix de ce chapitre, ajustée sur la première heure simulée avec cinq tranches et un horizon de cinq secondes, ajoute $-0.087$, $-0.007$, $0.022$, $0.091$ et
$0.220$ ticks au mid pour des déséquilibres compris dans $[0,0.2)$, $[0.2,0.4)$, \dots, $[0.8,1]$. Ici elle ajoute 0.091 et donne 99.59~: les données montrent que le mid pondéré
surestime ce que prédit le déséquilibre.
\end{example}

\section{Prix juste sur plusieurs plateformes}

Lorsqu'un instrument se négocie sur plusieurs plateformes, le carnet de chaque plateforme est une vue bruitée et retardée différemment d'un même prix efficient. Une plateforme dont la cotation
n'a pas changé depuis une seconde en dit moins qu'une cotation modifiée il y a une milliseconde~; une plateforme au carnet mince est plus bruitée que le marché primaire.

\begin{definition}[Prix juste consolidé]\label{def:hf:fair-value:consolidated}
Un \emph{prix juste consolidé}\index{prix juste consolidé} (consolidated fair price) est un \omterm{def:hf:fair-value:fair}{prix juste} qui combine les carnets de chaque plateforme sur laquelle un instrument
se négocie, chacun pondéré selon ce que ses prix apprennent du prix efficient, compte tenu de sa liquidité et de la date de sa dernière variation.
\end{definition}

Un \omterm{def:hf:fair-value:consolidated}{prix juste consolidé} n'est pas les meilleurs cours acheteur et vendeur nationaux (NBBO, livre~1, chapitre~9), qui est une construction réglementaire fondée sur les cotations protégées.
Le teneur de marché se soucie de l'information, et non de la protection~: une plateforme qui négocie rarement peut porter un meilleur cours périmé et mérite presque aucun poids.

La calibration la plus simple est une régression~: le mid futur de l'instrument, moins son mid courant, sur les écarts entre chaque source et le mid.
Sur le marché de ce chapitre (un suiveur B, coté sur deux plateformes, B et C), les poids ajustés sur la première heure valent 0.73 pour le microprix propre à B et 0.32 pour
celui de C~; elles totalisent environ un, comme il se doit pour deux mesures d'un même prix, et la seconde plateforme capte un tiers du poids.

\section{Prix juste sur plusieurs instruments}

\begin{definition}[Prix juste inter-instruments]\label{def:hf:fair-value:cross}
Un \emph{prix juste inter-instruments}\index{prix juste inter-instruments} (cross-instrument fair price) est un \omterm{def:hf:fair-value:fair}{prix juste} qui utilise aussi les prix des instruments liés, traduits
dans les unités propres à l'instrument (par un ratio de couverture, un ratio de conversion, une devise ou une base de fair value), de sorte qu'un mouvement d'un instrument lié
qui précède soit répercuté avant que le carnet propre à l'instrument ne le montre.
\end{definition}

Le teneur de marché d'un fonds indiciel coté (ETF) surveille le contrat à terme sur indice~; celui d'un certificat représentatif (chapitre~14) suit la cotation d'origine et la devise~;
un teneur de marché d'options suit le sous-jacent. L'instrument lié doit d'abord être mis dans les bonnes unités. Pour le contrat à terme d'un indice et d'un fonds
sur cet indice, la traduction est la base de fair value du livre~1, chapitre~21~; sur le marché de ce chapitre, les deux partagent un même prix efficient,
si bien qu'aucune traduction n'est nécessaire.

Une fois le mid du meneur ajouté à la régression, les poids deviennent 0.74 pour le microprix de B, 0.22 pour celui de C et 0.18 pour celui du meneur~; lorsque le meneur
bouge deux secondes avant B, le poids du meneur monte à 0.27, celui propre à B à 0.89, et celui de C tombe à $-0.09$. La régression exploite l'avance que le chapitre~8 va
estimer et exploiter.

\section{Filtrage~: le prix juste comme état}

Les régressions pondèrent les sources selon leur information moyenne. Un filtre pondère chaque observation à son arrivée, selon le temps écoulé depuis la
dernière observation.

\begin{definition}[Filtre de prix juste]\label{def:hf:fair-value:filter}
Un \emph{filtre de prix juste}\index{filtre de prix juste} (fair-price filter) traite le prix efficient comme un état caché, une marche aléatoire de variance $q$ par seconde, et
considère le prix de chaque source comme une observation de celui-ci, affectée d'une variance de bruit $r_j$. Pour une observation $y$ issue de la source $j$ à l'instant $t$, la mise à jour précédente
étant intervenue à $t'$, on propage et on met à jour~:
\[
\Sigma\leftarrow\Sigma+q\,(t-t'),\qquad K=\frac{\Sigma}{\Sigma+r_j},\qquad \hat P_t\leftarrow\hat P_t+K\,(y-\hat P_t),\qquad \Sigma\leftarrow(1-K)\,\Sigma .
\]
\end{definition}

C'est le filtre de Kalman à niveau local du livre~4, chapitre~19, appliqué à des temps irréguliers~; une période calme augmente $\Sigma$ et donc le poids de la
prochaine observation, ce qu'exige précisément une plateforme dont le cours est périmé (voir le \cref{lst:hf:fair-value:filter}).

\begin{proposition}[État stationnaire sous observations régulières]\label{prop:hf:fair-value:steady}
Si une source est observée toutes les $\Delta$ secondes, la variance a posteriori converge vers
\[
\bar\Sigma=\tfrac12\Bigl(-q\Delta+\sqrt{q^2\Delta^2+4q\Delta r}\Bigr),\qquad \bar K=\frac{\bar\Sigma+q\Delta}{\bar\Sigma+q\Delta+r}.
\]
\end{proposition}

\begin{proof}
Un point fixe de $\Sigma\mapsto (\Sigma+q\Delta)r/(\Sigma+q\Delta+r)$ vérifie $\Sigma^2+q\Delta\Sigma-q\Delta r=0$~; sa racine positive est
$\bar\Sigma$. L'application est croissante
et concave, de pente inférieure à un au point fixe, si bien que l'itération converge depuis tout point de départ positif.
\end{proof}

Avec $q=0.05$ ticks carrés par seconde, $r=0.1$ et dix observations par seconde, $\bar\Sigma=0.02$ et $\bar K=0.2$~: chaque nouveau mid déplace le \omterm{def:hf:fair-value:fair}{prix juste}
d'un cinquième de sa surprise. Le filtre du chapitre utilise $q=0.047$, le taux de variance du prix efficient sur la première heure, et pour chaque source la
variance de son écart au mid cinq secondes plus tard, 0.104, 0.123 et 0.134 pour B, C et le meneur.

\begin{omfigure}
\begin{tikzpicture}[font=\small,>=stealth,box/.style={draw,rounded corners,align=center,minimum height=1.0cm,text width=2.9cm,fill=omDef!6}]
\node[box] (b) at (0,2.2) {carnet de la plateforme B\\ mid, files};
\node[box] (c) at (0,0.8) {carnet de la plateforme C\\ mid, files};
\node[box] (a) at (0,-0.6) {meneur A\\ (le contrat à terme sur indice)};
\node[box,fill=omProp!10,text width=3.4cm] (f) at (5.2,0.8) {régression ou filtre~:\\ poids selon l'information};
\node[box,fill=omThm!8,text width=2.3cm] (p) at (9.6,0.8) {prix juste $\hat P_t$\\ de B};
\draw[->] (b.east) -- (f.west);
\draw[->] (c.east) -- (f.west);
\draw[->] (a.east) -- node[below right,font=\scriptsize,pos=0.35] {en unités de B} (f.west);
\draw[->] (f) -- (p);
\end{tikzpicture}

\omcaption{Le \omterm{def:hf:fair-value:fair}{prix juste} du chapitre pour l'instrument B~: le carnet propre de B sur sa plateforme principale, son carnet sur une deuxième plateforme et un instrument meneur
traduit en unités de B, combiné par des poids qui traduisent ce que chacun révèle sur le prix efficient de B.}\label{fig:hf:fair-value:sources}
\end{omfigure}

\section{Évaluer un prix juste}

Un \omterm{def:hf:fair-value:fair}{prix juste} se juge à ce qu'il prévoit. Dans une simulation, le prix efficient est connu et l'erreur peut donc être mesurée par rapport à lui~; sur un marché
réel, seuls les prix futurs sont observés, et la cible naturelle est le mid quelques secondes plus tard. Les deux cibles sont utilisées ci-dessous, sur une grille
de 0.1 seconde pendant la deuxième heure simulée, chaque paramètre ayant été ajusté sur la première heure.

\begin{center}
\small
\begin{tabular}{@{}lcccc@{}}
\toprule
& \multicolumn{2}{c}{RMSE (ticks), avance de 0.5 s} & \multicolumn{2}{c}{variation par rapport au mid (\%)}\\
estimateur du \omterm{def:hf:fair-value:fair}{prix juste} de B & prix efficient & mid dans 5 s & avance de 0.5 s & avance de 2 s\\
\midrule
mid & 0.962 & 0.549 & 0.0 & 0.0\\
mid pondéré & 0.963 & 0.542 & $-1.3$ & $+0.1$\\
microprix & 0.957 & 0.536 & $-2.4$ & $-1.7$\\
consolidé, régression & 0.947 & 0.542 & $-1.2$ & $-3.0$\\
consolidé, filtre & 0.961 & 0.559 & $+1.8$ & $-0.3$\\
inter-instruments, régression & 0.930 & 0.529 & $-3.5$ & $-12.8$\\
inter-instruments, filtre & 0.941 & 0.544 & $-0.8$ & $-8.7$\\
\bottomrule
\end{tabular}
\end{center}

Les deux dernières colonnes sont des variations de l'erreur par rapport au mid cinq secondes plus tard, moyennées sur trois marchés simulés. Trois enseignements se dégagent.
Le carnet de l'instrument lui-même apprend peu~: sur ce marché, tous les estimateurs construits à partir de lui sont à moins de 2.5~\% du mid. Le meneur en apprend davantage,
et d'autant plus qu'il est en avance (voir la \cref{fig:hf:fair-value:lags})~; avec une avance de deux secondes, le \omterm{def:hf:fair-value:cross}{prix juste inter-instruments} élimine un huitième de
l'erreur. Enfin, le filtre, dont les variances ont été fixées par une règle empirique, fait moins bien que la régression, qui a été ajustée à la cible même sur laquelle il est
évalué~: il faut calibrer un \omterm{def:hf:fair-value:fair}{prix juste} sur ce qu'il servira à prévoir.

\begin{omfigure}
\begin{tikzpicture}
\begin{axis}[width=11cm,height=5.2cm,xlabel={\small avance de A sur B (secondes)},ylabel={\small variation de RMSE (\%)},xmode=log,log basis x=10,
  xmin=0.08,xmax=2.5,ymin=-14,ymax=2,xtick={0.1,0.25,0.5,1,2},xticklabels={0.1,0.25,0.5,1,2},tick label style={font=\footnotesize},
  grid=major,grid style={gray!25},table/col sep=comma,legend style={font=\scriptsize,at={(0.5,-0.32)},anchor=north},legend columns=4]
\addplot[black,thick,mark=*,mark size=1.2pt] table[x=lag,y=micro]{figdata/hft/02-fair-value/lags.csv};
\addplot[omDef,thick,dashed,mark=square*,mark size=1.2pt] table[x=lag,y=consolidated]{figdata/hft/02-fair-value/lags.csv};
\addplot[omThm,thick,mark=triangle*,mark size=1.5pt] table[x=lag,y=cross]{figdata/hft/02-fair-value/lags.csv};
\addplot[omProp,thick,dotted,mark=diamond*,mark size=1.5pt] table[x=lag,y=cross_filter]{figdata/hft/02-fair-value/lags.csv};
\legend{microprix,consolidé,{inter-instruments, régression},{inter-instruments, filtre}}
\end{axis}
\end{tikzpicture}

\omcaption{Variation de l'erreur de chaque \omterm{def:hf:fair-value:fair}{prix juste}, par rapport à celle du mid, comme prévision du mid de B cinq secondes plus tard, selon le nombre de secondes dont
le meneur A se déplace avant B (échelle logarithmique)~; moyenne sur trois heures simulées chacune. Données~: \texttt{hf\_fairvalue.by\_lag}.}\label{fig:hf:fair-value:lags}
\end{omfigure}

La \Cref{fig:hf:fair-value:window} montre pourquoi. Après le plus grand mouvement de dix secondes de la deuxième heure, le prix efficient de B grimpe de cinq ticks, son mid suit
un tick à la fois pendant plusieurs secondes, et le \omterm{def:hf:fair-value:cross}{prix juste inter-instruments} passe devant le mid dès que le carnet du meneur a bougé. Dans la \cref{fig:hf:fair-value:lags}, avec trois marchés par point, les courbes sont bruitées (l'avance d'une demi-seconde montre par hasard moins que celle d'un quart de seconde)~; la tendance, elle, ne l'est pas.

\begin{omfigure}
\begin{tikzpicture}
\begin{axis}[width=11cm,height=5.2cm,xlabel={\small secondes depuis le mouvement},ylabel={\small prix (ticks)},tick label style={font=\footnotesize},
  xmin=-5,xmax=15,grid=major,grid style={gray!25},table/col sep=comma,scaled y ticks=false,
  yticklabel style={/pgf/number format/.cd,fixed,precision=0,1000 sep={}},
  legend style={font=\scriptsize,at={(0.5,-0.32)},anchor=north},legend columns=3]
\addplot[black,thick,const plot] table[x=t,y=truth]{figdata/hft/02-fair-value/window.csv};
\addplot[omDef,thick,const plot,dashed] table[x=t,y=mid]{figdata/hft/02-fair-value/window.csv};
\addplot[omThm,thick] table[x=t,y=cross]{figdata/hft/02-fair-value/window.csv};
\legend{prix efficient,mid,prix juste inter-instruments}
\end{axis}
\end{tikzpicture}

\omcaption{Instrument B autour du plus grand mouvement de dix secondes de son prix efficient dans la deuxième heure simulée (avance de A sur B de deux secondes)~: le prix efficient de B,
son mid et son \omterm{def:hf:fair-value:cross}{prix juste inter-instruments} sur une grille de 0.1 seconde. Données~: \texttt{hf\_fairvalue.window}.}\label{fig:hf:fair-value:window}
\end{omfigure}

\begin{remark}[Ce qui fait un bon résultat]\label{rem:hf:fair-value:good}
Une erreur inférieure de 3~\% ne semble pas grand-chose. Un teneur de marché gagne une fraction de la fourchette sur chaque exécution et perd la sélection adverse que son \omterm{def:hf:fair-value:fair}{prix juste}
n'a pas su anticiper (chapitre 1)~; le chapitre 6 montre que les exécutions qui jouent contre lui se concentrent aux moments où son \omterm{def:hf:fair-value:fair}{prix juste} était
le plus faux. La valeur d'un \omterm{def:hf:fair-value:fair}{prix juste} se mesure en sélection adverse évitée, par exécution, multipliée par toutes les exécutions d'une année.
\end{remark}

\section{Tutoriel~: trois marchés, sept prix justes}\label{tut:hf:fair-value:tut}

\begin{tutorial}
\textbf{Objectif.} Construire les \omterm{def:hf:fair-value:fair}{prix justes} du chapitre pour un instrument suiveur coté sur deux plateformes, avec un meneur, les ajuster sur une heure et les évaluer
sur l'heure suivante. \textbf{État final~:} le tableau ci-dessus et les \cref{fig:hf:fair-value:lags,fig:hf:fair-value:window}.
\begin{tutsteps}
\item \textbf{Marchés.} \texttt{hf\_fairvalue.markets(lag)} simule A, ainsi que B et C dont le prix efficient est celui de A \texttt{lag} secondes plus tard,
chacun avec son propre flux d'ordres.
\item \textbf{Le filtre.} \texttt{firm.fairprice.FairFilter} est la \cref{def:hf:fair-value:filter} en code~; ses jumeaux en C++20 et en Rust reproduisent les
mêmes nombres sur un jeu de test commun de 3\,000 observations.
\omcode{code/firm/fairprice/firm_fairprice.py}{68}{92}{Le filtre de prix juste~: propager pendant le temps écoulé, puis mettre à jour.}\label{lst:hf:fair-value:filter}
\item \textbf{Les régressions.} Le mid futur moins le mid sur les écarts entre chaque source et le mid, ajustés sur la première heure.
\omcode{code/hft/02-fair-value/python/hf_fairvalue.py}{100}{113}{Prix justes par régression~: plan d'expérience et pondérations.}\label{lst:hf:fair-value:weights}
\item \textbf{Score.} \texttt{table(lag)} moyenne les erreurs sur trois graines~; \texttt{fig\_fairvalue.py} écrit les données des figures.
\end{tutsteps}
\textbf{Pour aller plus loin.} Rendre la plateforme C peu profonde et lente et observer son poids baisser~; faire varier l'avance du meneur au cours de la journée et réajuster les poids
toutes les dix minutes.
\end{tutorial}

\section{Construction~: le moteur de prix juste}\label{bld:hf:fair-value:fairprice}

\begin{build}
\bfield{Rôle} \omterm{def:hf:fair-value:fair}{Prix justes} en flux pour chaque Quoter du carnet, avec un cœur assez rapide pour le hot path.
\bfield{Interface} \texttt{weighted\_mid}, \texttt{imbalance}, \texttt{fit\_micro(imb, spread, dmid, n)}, \texttt{micro(\dots, g)}~;
\texttt{FairFilter(q, r).update(t, src, y) -> estimate}, \texttt{predict(t)}~; \texttt{run\_filter}~; \texttt{rmse}.
\bfield{Règles} Les paramètres sont ajustés sur des données antérieures à la période où ils sont utilisés. Les sources sont traduites dans les unités de l'instrument
avant d'entrer. La première observation initialise l'état. Les filtres C++20 (\texttt{cpp/firm\_fairprice.hpp}) et Rust (\texttt{rust/}) effectuent les
mêmes opérations dans le même ordre.
\bfield{Tests d'acceptation} \texttt{code/firm/fairprice/tests/}~: le mid pondéré et le microprix à la main~; une étape de Kalman à la main~; le jeu de test
reproduit par la référence Python et, dans \texttt{cpp/} et \texttt{rust/}, par les jumeaux, à $10^{-9}$ tick près.
\bfield{Pour aller plus loin} Bruit dépendant de l'ancienneté $r_j(t-t_j)$~; un filtre à deux états ayant l'avance du meneur comme état~; réajustement en ligne des pondérations.
\end{build}

\begin{omsources}
\item S.~Stoikov, «~The micro-price: a high-frequency estimator of future prices~», \emph{Quantitative Finance} 18(12), 2018.
\item Le livre~4, chapitre~19 (le filtre de Kalman) et le livre~7, chapitre~8 (le microprix et les variables tirées du carnet d'ordres).
\end{omsources}

\section{Exercices}

\begin{exercise}[$\star$]\label{exo:hf:fair-value:1}
Un carnet affiche 300 actions demandées à 99 et 100 actions offertes à 100. Calculer le mid pondéré.
\end{exercise}

\begin{exercise}[$\star$]\label{exo:hf:fair-value:2}
Avec la table de microprix du chapitre, quel est le microprix d'un carnet à un tick 99/100 dont le déséquilibre vaut 0.9~?
\end{exercise}

\begin{exercise}[$\star$]\label{exo:hf:fair-value:3}
Le filtre dispose de $\hat P=100$ avec $\Sigma=0.02$~; $q=0.05$, $r=0.1$. Un mid de 101 arrive 0.4 seconde après la dernière mise à jour. Donner le gain et le nouveau
\omterm{def:hf:fair-value:fair}{prix juste}.
\end{exercise}

\begin{exercise}[$\star\star$]\label{exo:hf:fair-value:4}
En utilisant la \cref{prop:hf:fair-value:steady}, calculer $\bar\Sigma$ et $\bar K$ pour $q=0.05$, $r=0.1$ et dix observations par seconde, puis dire ce
qui arrive à $\bar K$ lorsque la source ralentit à une observation par seconde.
\end{exercise}

\begin{exercise}[$\star\star$]\label{exo:hf:fair-value:5}
Pourquoi un filtre dont les variances sont fixées au jugé peut-il prévoir moins bien que le mid simple, alors qu'une régression ajustée sur le même historique ne le fait pas~?
\end{exercise}

\begin{exercise}[$\star\star$]\label{exo:hf:fair-value:6}
À partir de la \cref{fig:hf:fair-value:lags}, lire l'amélioration de la régression inter-instruments pour une avance d'une seconde.
\end{exercise}

\begin{exercise}[$\star\star\star$]\label{exo:hf:fair-value:7}
\emph{Programmation.} Calculer \texttt{table(1.0)} et donner la variation de l'erreur du filtre inter-instruments par rapport au mid, cinq secondes à l'avance.
\end{exercise}

\begin{exercise}[$\star\star\star$]\label{exo:hf:fair-value:8}
\emph{Trouver la faille.} «~Notre nouvelle table de microprix réduit l'erreur de prévision de 9~\% sur les données d'aujourd'hui, sur lesquelles nous l'avons ajustée.~»
\end{exercise}

\section{Problème~: à qui appartient le prix}

\begin{problem}[{Problème du week-end --- à qui appartient le prix}]\label{pb:hf:fair-value:1}
Un fonds indiciel coté (ETF) se négocie sur deux plateformes et suit son contrat à terme sur indice. Il faut donner à son teneur de marché un seul nombre, mis à jour à chaque message.

\medskip\noindent\textbf{Partie I --- Un carnet.}
\begin{enumerate}
\item Définir le \omterm{def:hf:fair-value:fair}{prix juste} et le distinguer de la fair value d'un contrat à terme.
\item Écrire le mid pondéré et expliquer son sens.
\item Comment ajuste-t-on une table de microprix, et pourquoi seulement sur des carnets à un tick~?
\item Donner la table du chapitre et le microprix d'un carnet 99/100 à un déséquilibre de 0.75.
\item De combien le microprix améliore-t-il le mid comme prévision à cinq secondes~?
\end{enumerate}

\medskip\noindent\textbf{Partie II --- Plusieurs plateformes, plusieurs instruments.}
\begin{enumerate}[resume]
\item Définir un \omterm{def:hf:fair-value:consolidated}{prix juste consolidé} et dire pourquoi il ne se confond pas avec les meilleurs cours acheteur et vendeur nationaux (NBBO).
\item Donner les poids de la régression sur les microprix de B et de C, et interpréter leur somme.
\item Définir un \omterm{def:hf:fair-value:cross}{prix juste inter-instruments} et dire ce qui doit se produire sur le prix du meneur avant qu'il n'entre.
\item Donner les trois poids avec le meneur à une demi-seconde d'avance et à deux secondes d'avance.
\end{enumerate}

\medskip\noindent\textbf{Partie III --- Le filtre.}
\begin{enumerate}[resume]
\item Écrire les étapes de propagation et de mise à jour du \omterm{def:hf:fair-value:filter}{filtre de prix juste}.
\item Démontrer la \cref{prop:hf:fair-value:steady}.
\item Pourquoi une source en retard reçoit-elle davantage de poids à sa prochaine observation~?
\item Donner le $q$ et le $r$ du chapitre et dire comment chacun a été estimé.
\item Pourquoi le filtre a-t-il moins bien prévu que la régression~?
\end{enumerate}

\medskip\noindent\textbf{Partie IV --- Le verdict.}
\begin{enumerate}[resume]
\item Énoncer le \emph{résultat nommé}~: l'erreur de prévision à cinq secondes du mid, du microprix et du \omterm{def:hf:fair-value:cross}{prix juste inter-instruments} à
une demi-seconde d'avance, ainsi que la part de l'amélioration inter-instruments apportée par le meneur.
\item Comment l'amélioration évolue-t-elle avec l'avance~?
\item Pourquoi trois heures simulées sont-elles insuffisantes pour faire confiance à la forme de la \cref{fig:hf:fair-value:lags} point par point~?
\item Que montre la \cref{fig:hf:fair-value:window}~?
\item Comment transforme-t-on la qualité d'un \omterm{def:hf:fair-value:fair}{prix juste} en argent~?
\item En une phrase~: qu'est-ce qui rend un \omterm{def:hf:fair-value:fair}{prix juste} meilleur que le mid~?
\end{enumerate}
\end{problem}

\section{Questions d'entretien}

\begin{interviewq}[$\star$ \iqroles{trader}]\label{iq:hf:fair-value:1}
La file acheteuse est dix fois plus longue que la file vendeuse. Où se situe le \omterm{def:hf:fair-value:fair}{prix juste}, et pourquoi~?
\end{interviewq}

\begin{interviewq}[$\star\star$ \iqroles{chercheur}]\label{iq:hf:fair-value:2}
Comment évalueriez-vous un modèle de \omterm{def:hf:fair-value:fair}{prix juste} lorsque le prix efficient n'est jamais observé~?
\end{interviewq}

\begin{interviewq}[$\star\star$ \iqroles{chercheur}]\label{iq:hf:fair-value:3}
Dérivez le gain de Kalman en régime permanent d'une marche aléatoire observée avec un bruit toutes les $\Delta$ secondes.
\end{interviewq}

\begin{interviewq}[$\star\star$ \iqroles{développeur}]\label{iq:hf:fair-value:4}
Un service de \omterm{def:hf:fair-value:fair}{prix juste} consomme les flux de huit plateformes. Que doit transporter chaque mise à jour, et que faites-vous lorsqu'un flux se tait~?
\end{interviewq}

\begin{interviewq}[$\star\star$ \iqroles{trader}]\label{iq:hf:fair-value:5}
Le \omterm{def:hf:fair-value:fair}{prix juste} de votre fonds utilise le contrat à terme sur indice. Le marché de ce contrat ferme plus tôt aujourd'hui. Qu'est-ce qui change~?
\end{interviewq}

\begin{interviewq}[$\star\star\star$ \iqroles{chercheur}]\label{iq:hf:fair-value:6}
Les mids de deux plateformes mesurent un même prix avec des bruits indépendants de variances $r_1$ et $r_2$. Quels poids minimisent la variance de la combinaison,
et quand est-il optimal d'ignorer la seconde plateforme~?
\end{interviewq}
